\section{논의}\label{sec:discussion}
수정된 근거는 제어기 구조와 튜닝을 구분한다. 기존 LQI는 공칭 추종과 부하 제거에서 우세하고, 동일 예산 튜닝에서는 매우 높은 피드백 권한을 택한 LQR이 공칭 비교에서 우세하다. 이 이점은 조건부이다. 센싱 지연은 지속적인 포화를 유발하고 공칭 이점을 제거한다. 적분 동작은 실현 가능한 일정 열부하를 안정적으로 제거하지만 LQI가 잡음, 포화, 누락 동특성에서 항상 최선은 아니다.

안티와인드업 수정은 과학적으로 중요하다. 수정 전에는 공정이 더 좁은 한계를 사용하는 동안 PID와 LQI가 공칭 한계에 대해 적분했다. 수정 결과는 제한 액추에이터에서 기존 LQR이 가장 좋다는 이전 가설을 지지하지 않는다. 재튜닝 이득에서는 LQR의 제한 IAE가 근소하게 가장 작지만 차이는 실용적으로 작고 모든 제어기에서 포화가 지배적이다.

같은 후보 수와 시나리오는 공정성을 높이지만 최적화 난이도를 같게 하지는 않는다. LQR은 두 가중치, PID와 LQI는 세 매개변수를 탐색하고 매개변수-응답 사상도 다르다. 따라서 결과는 통제된 유한 예산 비교로 보는 것이 적절하다. 독립 가우시안 실현의 통계적 신뢰도 역시 몬테카를로 정밀도를 나타낼 뿐 실제 센서 물리나 공정 매개변수 불확실성을 뜻하지 않는다.

화학적 역할은 의도적으로 제한했다. 모델은 비반응성 액체의 용액 열용량, 열전달, 열평형을 표현한다. 반응기 동역학, 농도, 산업 장비, 실험실 검증을 주장하지 않는다.
